\documentclass[tikz,border=5mm]{standalone} %% made with Claude.
% \usepackage{hsbcommon} %% changes the colours!
\usetikzlibrary{fadings,calc}

\begin{document}
% Fading: transparent in the middle, increasingly opaque toward the rim
\pgfdeclareradialshading{rimshade}{\pgfpoint{0bp}{0bp}}{%
  color(0bp)=(pgftransparent!100);
  color(24bp)=(pgftransparent!80);
  color(40bp)=(pgftransparent!35);
  color(47bp)=(pgftransparent!5);
  color(50bp)=(pgftransparent!0)}
\pgfdeclarefading{rim}{\pgfuseshading{rimshade}}

% Iridescent thin-film colours
\pgfdeclarehorizontalshading{iridescent}{100bp}{%
  color(0bp)=(cyan);
  color(22bp)=(magenta!80);
  color(40bp)=(yellow!70!orange);
  color(58bp)=(green!70!cyan);
  color(76bp)=(violet!80);
  color(100bp)=(cyan!60)}

% ---------- "balloon" letters A and I ----------
% Every stroke of a letter is an inflated "sausage": narrow (pinched) at
% its ends and bulging outwards in the middle, like a twisted balloon.
%   \bseg{P}{Q}{bulge}  -> closed outline of one balloon segment P->Q
%   \bglint{P}{Q}{bulge}-> curved highlight along the segment's upper side
\def\pinch{0.27cm}   % half-width at the pinched ends
\def\bcap{0.36cm}    % how far the rounded ends stick out
\newcommand{\bseg}[3]{%
  ($(#1)!\pinch!90:(#2)$)
  .. controls ($($(#1)!0.25!(#2)$)!#3!90:(#2)$)
          and ($($(#1)!0.75!(#2)$)!#3!90:(#2)$) ..
  ($(#2)!\pinch!-90:(#1)$)
  .. controls ($($(#2)!\pinch!-90:(#1)$)!-\bcap!($(#1)!\pinch!90:(#2)$)$)
          and ($($(#2)!\pinch!90:(#1)$)!-\bcap!($(#1)!\pinch!-90:(#2)$)$) ..
  ($(#2)!\pinch!90:(#1)$)
  .. controls ($($(#1)!0.75!(#2)$)!#3!-90:(#2)$)
          and ($($(#1)!0.25!(#2)$)!#3!-90:(#2)$) ..
  ($(#1)!\pinch!-90:(#2)$)
  .. controls ($($(#1)!\pinch!-90:(#2)$)!-\bcap!($(#2)!\pinch!90:(#1)$)$)
          and ($($(#1)!\pinch!90:(#2)$)!-\bcap!($(#2)!\pinch!-90:(#1)$)$) ..
  cycle}
\newcommand{\bglint}[3]{%
  ($($(#1)!0.22!(#2)$)!0.40*#3!90:(#2)$)
  .. controls ($($(#1)!0.40!(#2)$)!0.62*#3!90:(#2)$) ..
  ($($(#1)!0.62!(#2)$)!0.55*#3!90:(#2)$)}

% The strokes: P -> Q is oriented so that the "left" side faces up/left
% (towards the light), which is where the highlights go.
\newcommand{\AIletters}[1]{% #1 = bulge scale (1 = normal)
  \bseg{-1.3,-1.1}{-0.45,1.1}{#1*0.52cm}     % A: left leg
  \bseg{0.4,-1.1}{-0.45,1.1}{#1*0.52cm}      % A: right leg
  \bseg{-0.92,-0.3}{0.02,-0.3}{#1*0.36cm}    % A: crossbar
  \bseg{1.3,-1.1}{1.3,1.1}{#1*0.52cm}}       % I
\newcommand{\AIglints}{%
  \bglint{-1.3,-1.1}{-0.45,1.1}{0.52cm}
  \bglint{0.4,-1.1}{-0.45,1.1}{0.52cm}
  \bglint{-0.92,-0.3}{0.02,-0.3}{0.36cm}
  \bglint{1.3,-1.1}{1.3,1.1}{0.52cm}}

\begin{tikzpicture}
  \def\R{3cm} % radius -> 6 cm diameter

  % soft background so the transparent bubble is visible
  \shade[top color=blue!25, bottom color=blue!5]
    (-4,-4) rectangle (4,4);

  % balloon letters "AI", drawn first so the bubble film lies over them
  % shift that moves the letters' area centroid (0.165, -0.083)
  % onto the bubble centre (0,0)
  \begin{scope}[shift={(-0.165,0.083)}]
    \path[draw=red!45!black, line width=1.2mm, line join=round] \AIletters{1}; % rim
    \fill[red!85!orange] \AIletters{1};                                       % body
    \begin{scope}[shift={(-0.05,0.05)}]                                    % soft sheen
      \def\pinch{0.12cm}
      \fill[red!25!white, opacity=0.35] \AIletters{0.6};
    \end{scope}
    \draw[white, opacity=0.9, line width=0.9mm, line cap=round] \AIglints;  % glints
  \end{scope}

  % very faint body tint
  \fill[cyan!20, opacity=0.25] (0,0) circle (\R);

  % iridescent film, strongest at the rim
  \shade[shading=iridescent, shading angle=35,
         path fading=rim, fit fading=true]
    (0,0) circle (\R);
  % second layer, rotated, for swirling colours
  \shade[shading=iridescent, shading angle=-60, opacity=0.5,
         path fading=rim, fit fading=true]
    (0,0) circle (\R);

  % thin outline
  \draw[white, opacity=0.6, line width=0.5pt] (0,0) circle (\R);

  % main reflection: curved "window" in the upper left
  \fill[white, opacity=0.85]
    (128:2.35) arc (128:158:2.35)
    -- (158:1.75) arc (158:128:1.75) -- cycle;
  % long glossy arc along the upper rim
  \draw[white, opacity=0.7, line width=2.5pt, line cap=round]
    (100:2.72) arc (100:165:2.72);
  % small bright spot
  \fill[white, opacity=0.9] (-0.75,1.95) circle (0.12);
  % secondary reflection, lower right
  \draw[white, opacity=0.5, line width=1.5pt, line cap=round]
    (-55:2.65) arc (-55:-20:2.65);
\end{tikzpicture}
\end{document}